\subsection{Sentence-Level Hallucination Detection}
\label{sec:sentence_level}

\paragraph{Comparing detection methods at the sentence level.}
Table~\ref{tab:main} reports PRR for each baseline across
RAGTruth \citep{niu2024ragtruth}, FactScore \citep{min2023}, and
LongFact \citep{wei2024}, evaluated under the three instruction-tuned
LLMs introduced in Section~\ref{sec:experiments}. The Temp-ViT
(Stage~3) row corresponds to our per-atom output trained under the
passage-level refinement of Section~\ref{sec:stage3}, while
Temp-ViT (Stage~2) corresponds to the per-window MIL output of
Section~\ref{sec:temporal_head} aggregated per claim by averaging its
in-span scores. Together they answer (Q\textsubscript{1}): the passage
refinement step is what gives the per-atom scores room to differentiate
uncertain claims from confident ones within the same response, where
the consistency-based and information-based families compress all
claims of a response to a single sequence-level rank. Extended results
on the corresponding AUROC, ECE, and Brier metrics are deferred to
Appendix~\ref{app:extended}.

\begin{table*}[t]
\centering
\caption{\textbf{Long-form hallucination detection --- token level.} We
report AUROC per (LLM, dataset) cell. Models: Mistral-7B-Instruct-v0.2
(Mis-7B-Inst), Llama-3.1-8B-Instruct (LlaMa-8B-Inst),
Llama-3.1-70B-Instruct (LlaMa-70B-Inst). Datasets: RAGTruth
\citep{niu2024ragtruth}, FactScore \citep{min2023}, LongFact
\citep{wei2024}. Per-token gold mask built from atom spans via
longer-span-wins on containment, earlier-span-wins on partial overlap;
uncovered tokens dropped. Best score per column in \textbf{bold},
second-best \underline{underlined}. $^\dagger$ marks cells whose data
comes from a closely related model version
(Mistral-7B-Instruct-v0.1 for RAGTruth, Llama-3-8B-Instruct for
LongFact). Cells with ``---'' indicate experiments not yet run.}
\label{tab:main}
\setlength{\tabcolsep}{2.5pt}
\renewcommand{\arraystretch}{1.05}
\scriptsize
\adjustbox{max width=\textwidth}{%
\begin{tabular}{l ccc ccc ccc}
\toprule
 & \multicolumn{3}{c}{Mis-7B-Inst} & \multicolumn{3}{c}{LlaMa-8B-Inst} & \multicolumn{3}{c}{LlaMa-70B-Inst} \\
\cmidrule(lr){2-4} \cmidrule(lr){5-7} \cmidrule(lr){8-10}
Method & RAGTruth & FactScore & LongFact & RAGTruth & FactScore & LongFact & RAGTruth & FactScore & LongFact \\
\midrule
SemanticEntropy           & --- & --- & --- & --- & --- & --- & --- & --- & --- \\
DegMat (NLI)              & --- & --- & --- & --- & --- & --- & --- & --- & --- \\
EigValLaplacian (NLI)     & --- & --- & --- & --- & --- & --- & --- & --- & --- \\
SAR                       & --- & --- & --- & --- & --- & --- & --- & --- & --- \\
LUQ                       & --- & --- & --- & --- & --- & --- & --- & --- & --- \\
KernelLanguageEntropy     & --- & --- & --- & --- & --- & --- & --- & --- & --- \\
\midrule
Perplexity                & --- & --- & --- & --- & --- & --- & --- & --- & --- \\
MeanTokenEntropy          & --- & --- & --- & --- & --- & --- & --- & --- & --- \\
MCNSE                     & --- & --- & --- & --- & --- & --- & --- & --- & --- \\
TokenSAR                  & --- & --- & --- & --- & --- & --- & --- & --- & --- \\
\midrule
CocoaMSP                  & --- & --- & --- & --- & --- & --- & --- & --- & --- \\
CocoaPPL                  & --- & --- & --- & --- & --- & --- & --- & --- & --- \\
CocoaMTE                  & --- & --- & --- & --- & --- & --- & --- & --- & --- \\
\midrule
Act-ViT (F-v2, N\textsubscript{eff}=6)
                          & \underline{0.570}$^{\dagger}$ & \underline{0.676} & \underline{0.638} & --- & \textbf{0.673} & \underline{0.578}$^{\dagger}$ & --- & --- & --- \\
Temp-ViT (Stage 3, ours)  & \textbf{0.865}$^{\dagger}$ & \textbf{0.696} & \textbf{0.692} & --- & \underline{0.653} & \textbf{0.601}$^{\dagger}$ & --- & --- & --- \\
\bottomrule
\end{tabular}%
}
\end{table*}


\paragraph{Run-time and performance trade-offs.} Temp-ViT scores every
generated token in a single forward pass through the frozen encoder and
the Mamba-2 temporal head, so the cost of obtaining a sentence-level
hallucination probability is amortized across the whole passage rather
than scaling with the number of claims. We quantify this in
Figure~\ref{fig:inference_time} of Section~\ref{sec:non_static}, which
shows roughly an order-of-magnitude wall-clock advantage over
sequence-level activation-tensor baselines that must be re-invoked per
claim. The training cost for the passage-level refinement stage is
under three hours on a single H200 across all three LLMs; the encoder
is trained once and reused across stages.

We also study the effect of more aggressive activation pooling on this
trade-off. Following the same protocol as Section~\ref{sec:experiments},
we sweep the pooling target on LongFact and report Test AUC vs.\
training runtime in Figure~\ref{fig:ablation}. Aggressive pooling
($(L_p, N_p) = (4, 20)$) preserves most of the per-atom AUROC while
shrinking training to a few minutes, in line with the
scalability claims of (Q\textsubscript{3}).

\begin{figure}[t]
\centering
\fbox{\begin{minipage}{0.95\linewidth}\centering\vspace{4em}\textit{[Figure 3: Ablation study on the pooling hyperparams $(L_p, N_p)$ --- placeholder]}\vspace{4em}\end{minipage}}
\caption{Ablation study on the pooling hyperparams $(L_p, N_p)$. Probe[$*$] AUC is indicated by dashed green line in the left plot.}
\label{fig:ablation}
\end{figure}

